% from ega0-22-fr.tex
\oldpage[0\textsubscript{IV-1}]{203}

\begin{env}[22.5.5]
\label{0.22.5.5-fr}
Soient $k$ un corps de caractéristique $p>0$, $A$ un anneau local, $k'$ une
extension \emph{radicielle finie} de $k$ telle que $k'^p\subset k$,
$\mathfrak m$ l'idéal maximal de $A$, $K=A/\mathfrak m$ son corps résiduel~;
rappelons que $A'=A\otimes_kk'$ est un anneau local dont le corps résiduel est
le même que celui de l'anneau local $K\otimes_kk'$ (Bourbaki,
\emph{Alg. comm.}, chap.~V, §~2, n\textsuperscript{o}~3, lemme~4). Notons que
l'on a $A'_{(K)}=A'\otimes_AK=K\otimes_kk'$, et par suite (21.3.6)
\begin{equation}
\label{0.22.5.5.1-fr}
  \Theta_{A'_{(K)}/K}=\Theta_{k'/k}\otimes_kK
  \tag{22.5.5.1}
\end{equation}
à un isomorphisme canonique près. En outre, on sait (21.4.7) que l'on a
$\Xi_{k'/k}=0$, autrement dit l'homomorphisme canonique (21.3.2.2)
\[
  \pi_{k'/k}:\Theta_{k'/k}\longrightarrow\Omega^1_k
\]
est injectif, si bien que l'on a une \emph{injection canonique}
\begin{equation}
\label{0.22.5.5.2-fr}
  \Theta_{A'_{(K)}/K}\longrightarrow\Omega^1_k\otimes_kK
  \tag{22.5.5.2}
\end{equation}
qui s'explicite de la façon suivante (compte tenu de (22.5.5.1))~: pour tout
$x'\in k'$, à l'élément
$d_{A'_{(K)}/K}(x'\otimes1_K)\otimes1_K$ de
$\Omega^1_{A'_{(K)}/K}\otimes_{A'_{(K)}}K$, on fait correspondre l'élément
$d_k(x'^p)\otimes1_K$ de $\Omega^1_k\otimes_kK$.
\end{env}

\begin{lemma}[22.5.6]
\label{0.22.5.6-fr}
Avec les notations de (22.5.5)~:
\begin{enumerate}
  \item[\rm(i)] Le noyau de l'homomorphisme canonique
  $\pi_{A'_{(K)}/K}$ (21.3.2.2) s'identifie par (22.5.5.2) à
  $(\Theta_{k'/k}\otimes_kK)\cap\Upsilon_{K/k}$.
  \item[\rm(ii)] Soient $A_1=A/\mathfrak m^2$,
  $A'_1=A_1\otimes_kk'=A'/\mathfrak m^2A'$~; alors l'homomorphisme
  caractéristique $\chi'_{A'_1/A_1}$ (22.4.6.2) s'identifie à la restriction de
  l'homomorphisme caractéristique
  $\chi_{A/k}:\Upsilon_{K/k}\longrightarrow V=\mathfrak m/\mathfrak m^2$
  (20.6.24) à $(\Theta_{k'/k}\otimes_kK)\cap\Upsilon_{K/k}$.
\end{enumerate}
\end{lemma}

\begin{proof}
(i) En vertu de (20.6.9), appliqué en remplaçant $A$ par le corps premier, $B$
par $k$, $C$ par $K$ et $E$ par $A_1=A/\mathfrak m^2$ (extension de $K$ par
$V=\mathfrak m/\mathfrak m^2$), il suffit (cf. (22.5.5)) de vérifier que si
$x'\in k'$ et si $z\in A_1$ est un élément dont la classe mod.
$\mathfrak m/\mathfrak m^2$ est $x'^p\in k$, alors l'image de
$d_k(x'^p)\otimes1_K$ par l'homomorphisme canonique
$\Omega^1_k\otimes_kK\longrightarrow\Omega^1_{A_1}\otimes_{A_1}K$ est
$d_{A_1}z\otimes1_K$, ce qui résulte des définitions.

(ii) Comme $A'_{(K)}=K\otimes_kk'=(A'_1)_{(K)}$ (en vertu de la relation
$K=A_1/(\mathfrak m/\mathfrak m^2)$), $\Theta_{A'_{(K)}/K}$ s'identifie à
$\Theta_{(A'_1)_{(K)}/K}$ et la vérification résulte du même calcul que dans
(i) et de la définition (22.4.5.2) de $\pi'_{(A'_1)_{(K)}/K}$.
\end{proof}

\begin{proposition}[22.5.7]
\label{0.22.5.7-fr}
Soient $k$ un corps de caractéristique $p>0$, $A$ une $k$-algèbre qui est un
anneau local régulier, $\mathfrak m$ son idéal maximal, $K=A/\mathfrak m$ son
corps résiduel~; soit $V=\mathfrak m/\mathfrak m^2$, considéré comme
$K$-espace vectoriel. Soient $k'$ une extension finie de $k$ telle que
$k'^p\subset k$, et $A'=A\otimes_kk'$~; pour que l'anneau local $A'$ soit
régulier, il faut et il suffit que la restriction de l'homomorphisme
caractéristique $\chi_{A/k}:\Upsilon_{K/k}\longrightarrow V$ (20.6.24) à
$(\Theta_{k'/k}\otimes_kK)\cap\Upsilon_{K/k}$ soit injective.
\end{proposition}

\begin{proof}
Il s'agit de prouver que la condition de l'énoncé est équivalente à
(22.5.1.1), en désignant par $\mathfrak m'$ l'idéal maximal de $A'$, par
$K'=A'/\mathfrak m'$ son corps résiduel. Avec les notations de la remarque
(22.5.2, (i)), on a $A'_1=A_1\otimes_kk'$. Or, soit
$(a_i)_{1\leqslant i\leqslant r}$ une $p$-base de $k'$ sur $k$, de sorte que
$a_i^p=b_i\in k$, et $k'$ est isomorphe au quotient de l'anneau de

